\documentclass[leqno, a4paper, 12pt]{jsarticle}
\usepackage{amssymb}
\usepackage{amsthm}
\usepackage{amsmath}
\usepackage{comment}
\usepackage{url}

\usepackage{enumitem}
%\usepackage{showkeys}


%定理環境 thmはあまり使わずpropをなるべく使う。
%主張は基本的に証明の中で使い、そのため番号を付けない。
%注意環境を付けてみた。
\newtheorem{thm}{定理}
\newtheorem{prop}[thm]{命題}
\newtheorem{cor}[thm]{系}
\newtheorem{lem}[thm]{補題}
\newtheorem{df}[thm]{定義}
\newtheorem{exam}[thm]{例}
\newtheorem{fact}[thm]{事実}
\newtheorem*{claim}{主張}
\newtheorem*{cau}{注意}
\newtheorem*{rmk}{注意}
\renewcommand{\proofname}{証明}
%矢印たち。
%\toは\rightarrowと同じ。\getsは\leftarrowと同じ。同値は\iff
\newcommand{\To}{\Longrightarrow}
\newcommand{\Gets}{\LongLeftarrow}
%黒板太字
\newcommand{\nn}{\mathbb{N}}
\newcommand{\zz}{\mathbb{Z}}
\newcommand{\qq}{\mathbb{Q}}
\newcommand{\pp}{\mathbb{P}}
\newcommand{\rr}{\mathbb{R}}
\newcommand{\cc}{\mathbb{C}}
%無限大
\newcommand{\mgn}{\infty}
%差集合は\setminusで統一する。等号の否定は\neq
%真部分集合は\subsetneqq　偏微分は\partial
%ディスプレイスタイル。\dfracでディスプレイ分数
\newcommand{\dpst}{\displaystyle}

\title{論文メモ
\large{\\--- 位相的特徴付け ---}
}
\author{ナンブ　キトラ}
\date{\today}
\begin{document}
\maketitle

本棚を整理したいので，
溜まっている論文のメモを書いて未来への手紙とすることにする．

この論文の束は
特別な位相空間の特徴付けに関する
ものだと思う．

\section{本文}

\subsection{実直線とかの特徴付け}
論文
\cite{ward1936topological}
が最初に直線を特徴付けした論文のように見える．
この論文では直線$\rr$を
連結，局所連結，可分距離空間であって，
どの一点の補集合も二つの連結成分を持つような空間として
特徴付けしている．
Mooreの1920年の論文
\cite{MR1501148}
もこの話題に関する重要な論文のような気がするが，
書かれている数学語が古語なのでちょっと読みづらくて
わかんなかった．

論文
\cite{MR0094774}
で
M. E. Rudin 
は
実数直線を可分距離空間$E$でなおかつ次の二つの条件1. 
その連結成分が一点かarcになりそしてarcになる成分$A$の
どんな内点も$E\setminus A$の極限点にならない;
2.$E$のどんな近傍も境界集合が有限集合になるような基本近傍系をもつ，
という空間として特徴付けした．

Bing の論文
\cite{MR1530892}
は学術論文というよりかは実数の特徴付けに関する
解説論文である．
その中で実数直線$L$
は四つの公理; 公理１：$L$ is a linear topological space 
(この条件は今ふうに言うとLOTSであるということである); 
公理２： $L$は最大値と最小値を持たない; 
公理３： $L$は連結; 
公理４ ：$L$は可分である，
で特徴付けできることがまず紹介されている．
そして公理4を
公理 4' ：$L$の互いに交わらない開集合の族はたかだか可算族である，
に置き換えたとき，
$L$が実数直線と同相であるかどうかは未解決問題であることを述べている．
この公理1,2, 3, 4'を満たす直線で
$\rr$と\textbf{同相ではない}ものは
\textbf{ススリン線}と呼ばれる．
このススリン線が存在するのかどうかと言う問題は
\textbf{ススリンの問題}と呼ばれている．
ここで注意して欲しいのはこの論文が発表されたのは1960年であることで，次いでススリン線にまつわる集合論的見識が深まったのは60年代の後半だと言うことである．ちなみに\textbf{強制法}の発明は1963年くらいらしい．
今ではススリンの問題はZFCと独立であると言うことが知られている．
こういう「歴史の痕跡」のような文献を見つけると楽しくなります．

論文
\cite{MR0257981}
は一般化対角線を引っこ抜いた直積空間が連結であることが
直線と円周の特徴付けになっているらしいが，
よくわからなかった．


論文
\cite{MR0282341}
も実数の位相的特徴付けの話だと思う．

論文
\cite{MR0866194}
では
位相群としての
$\rr$を特徴付けしている．
$G$を離散的ではない局所コンパクトな位相群としたとき，
$G$の真部分群であって閉であるものが常に離散的であることと$G$が1次元トーラスか$\rr$に位相群として同型であることは同値であることを証明している．


論文
\cite{MR1617361}
では，弧状連結なハウスドルフ空間が連続選択子を持つならば
それは一点か，$[0, 1)$, 
$[0, 1]$，もしくは$\rr$に同相であることを証明している．
ここで位相空間$X$が連続選択子を持つとは，
$\mathcal{F}(X)$を$X$の空でない閉集合全体にVietoris位相を入れたものとして，連続関数
$f\colon \mathcal{F}(X)\to X$であって
$f(F)\in F$を満たすものが存在することである．

論文
\cite{MR0283767}
ではorder topologyを$T_{1}$
空間であって二つもminimalな$T_{0}$
位相のupper boundとして表せるものとして特徴付けをした上で，$\rr$
をconnected separable な$T_{1}$
空間でその位相は二つのminimal
$T_{0}$位相のupper bound
になっているものとして特徴付けしている．



プレプリント
\cite{lipham2019compactification}
では
cut-point space
というものとコンパクト化の関係を見ているらしいが
よくわからず．

Brownの
論文
\cite{MR0126835}
はなんていうかなんか大事なことをしているような気がする．
$X$の部分集合列$\{V_{i}\}_{i\in \zz_{\ge 0}}$
に関して
$V_{i}$が$\rr^{n}$
に同相で$V_{i}\subset V_{i+1}$
で
$X=\bigcup_{i}V_{i}$
ならば$X$は
$\rr^{n}$に同相であることを証明している．


論文
\cite{MR0220251}
では局所連結な局所コンパクト距離空間が$\rr$
の連続単写像になるならばそれは五種類の空間に分類できることを証明している．
$\rr$, 
$8$, 
$\bigcirc\! \! -\! \!\bigcirc$, 
$\theta$, 
$\bigcirc\!-$, 
の五種類である．



\subsection{順序数}
論文
\cite{MR2891422}
では位相空間$X$の
点を分離するような$X$の集合族から
誘導される順序を用いて順序数を特徴付けしている．

\subsection{シェルピンスキーの三角形}
Vejnarの論文
\cite{MR2879369}
ではシェルピンスキーの三角形を位相的に特徴付けしている．
詳細を述べるのはちょっと難しいので論文を読まれたし．

\subsection{$0$次元空間}

記事
\cite{van2003topological}
や
記事
\cite{isolated2005}
にも色々書いてある．
論文
\cite{MR0657649}
にも以下の基本的な空間の特徴付けのまとめや表とかが書いてある．読むといいと思う．
この論文自体は以下に紹介する基本的な0次元空間の
完全写像による像に関する
論文である．

\subsubsection{カントール集合の特徴付け}
カントール集合
$C$は
ブラウワーの論文
\cite{brouwer1910structure}
で完全不連結性で孤立点を
持たないコンパクト距離化可能空間として
特徴付けられているらしいが
なんかとても読みにくい文章で書かれているのでそこまで自信がない．
多分Theorem5,と6がそれに該当すると思う．
100年前の英語の言葉遣いなのでなんか読みにくい．
ケクリス
\cite{MR1321597}
にも証明が載っている．

論文
\cite{MR0377836}
ではカントール集合を
コンパクト距離空間$X$であって，
１．$X$はコンパクトな開集合を少なくとも一つ持ち，そして非コンパクトな開集合も少なくとも一つもつ．
２．コンパクトな開集合は全て同相
3. 非コンパクトな開集合は全て同相
の性質を持つものとして特徴付けた．

論文
\cite{MR0139142}
では，
upper semi-continuous map
の話を用いて以下のことを証明されている．
可分距離空間$X$
が0次元で
孤立点を持たず，
ポーランドであるならば
連続全単射$f\colon X\to C$
が存在することを証明した．
全単射であるだけで同相ではないし，
同相にならない$X$も構成できる．
特に$X$がコンパクトの場合にはコンパクトからハウスドルフへの連続全単射なので同相になる．

\subsubsection{無理数の空間(ベール空間)の特徴付け}
無理数の空間$\pp=\rr \setminus \qq$の特徴付けは
論文
\cite{MR1512393}
でウリゾーンによって与えられた．
ちなみにこの空間
$\pp$
は可算離散空間の可算直積
$\nn^{\nn}$と同相であり，
この二つの空間の同相を与える典型的な写像は
\textbf{連分数}である．
位相的特徴付けとしては，
$0$次元ポーランド空間であって，
その任意のコンパクト部分集合は絶対に内点を持たない空間として特徴付けされる．
後半の条件から孤立点を持たないことがわかる．
ケクリス
\cite{MR1321597}
にも証明が載っている．




\subsubsection{有理数の特徴付け}

有理数の空間の特徴付けは
\cite{sierpinski1920propriete}
で初めて与えられたらしい．
フランス語で書かれているのであんまり読めない．
曰く，有理数の空間$\qq$とは
可算な距離空間であって，孤立点を持たない空間として特徴付けできる．なので孤立点がない可分距離空間の可算な
稠密部分集合を考えるとそれらは全て同相であり，$\qq$
に同相である．$\qq^{2}$とか$\qq^{n}$とかも全部
$\qq$に同相である．
なぜか知らないがこの特徴付けの簡単な証明が続々発見されている．詳細は
\cite{MR3949184}, 
\cite{MR4079948}, 
や
\cite{MR4217759}
を参照のこと．
\cite{MR4217759}が一番新しい．
また
プレプリント
\cite{francis2012two}
でもブラウワーの定理とシェルピンスキーの定理
が紹介されている．
論文
\cite{MR0644644}
ではカントール集合の特徴付けを用いて
$\qq$の特徴付けを証明している．

\subsubsection{$\qq\times C$}
van Mill
の論文
\cite{MR0597877}
によれば
アレキサンダーとウリゾーンの論文
\cite{MR1512393}
では$\qq\times C$
を
$0$次元な可分距離空間で$\sigma$コンパクトでコンパクト部分集合が常にない点を持たずどの開集合も可算ではないような
空間として特徴付けしていると言う．
しかしドイツ語が読めないので，
論文\cite{MR1512393}の記述は
私はよくわかっていない．
論文
\cite{MR1391294}
も参照のこと．


\subsubsection{$\qq\times \pp$}
同じく
van Millの論文
\cite{MR0597877}
では$\qq\times \pp$
を$0$次元可分距離空間であって，
可算個のポーランド部分空間で被覆され($\sigma$-complete)，どの開集合も完備ではなく$\sigma$-コンパクトではないような空間として特徴付けしている．


\subsection{任意の空でないコンパクト
距離空間がカントール集合の連続像になること}
趣がちょっと異なるが，
任意の空でないコンパクト距離空間がカントール
集合の連続像になるという定理に関する論文を紹介する．
ちなみにこの定理はアレキサンドロフとハウスドルフにより
独立に証明されたらしい．

論文
\cite{koudela2007hausdorff}
はサーベイだと思うが，掲載されているジャーナル？がよくわからない感じ．
まあサーベイである．
ハーン・マズルキーウィッチの定理，
つまり単位閉区間の連続像になるような空間の特徴づけも紹介されている．

論文
\cite{MR1650866}
では上記の定理の応用を紹介している．
セクション１では
space-filling curveを構成している．
セクション２ではなんか凸解析っぽい普遍空間の存在を述べる
Grazaslewiczの定理を証明している．
セクション３では$C[0, 1]$
に可分距離空間を等長に埋め込めるというBanach-Mazurの定理を証明している．
セクション４では連続関数であって，全ての有界関数を補完するものを構成している．詳しく述べると以下である．
\begin{thm}
連続関数$f\colon \rr\to \rr$
であって，以下の条件を満たすものが存在する：
両側に伸びた数列$\{y_{i}\}_{i\in \zz}$であって
$|y_{i}|\le 1$も満たす任意のものに対して
$t\in \rr$が存在して 
$y_{i}=f(t+i)$
を任意の$i\in \zz$
でみたす．
\end{thm}
セクション６では上の定理の変種を探索している．
セクション７では非常にゆっくりと$0$に収束する関数列を構成している．これのオリジナルはW. Rudinらしい
(A very slowly convergent sequence of continuous functions, Proc. Amer. Math. Soc., 39 (1973), 647--648)．

%READ!
%myplain is the same to amsplain style. 
\nocite{*}
\bibliographystyle{myplaindoidoi}
\bibliography{chch.bib}



\end{document}